% Part: first-order-logic
% Chapter: syntax-and-semantics
% Section: satisfaction

\documentclass[../../../include/open-logic-section]{subfiles}

\begin{document}

\olfileid[hi]{fol}{syn}{sat}

\olsection{\article{structure} \printtoken{S}{structure} में
  \article{formula} \printtoken{S}{formula} की संतुष्टि}

\begin{explain}
एक ओर पदों और !!{formula}s जैसी अभिव्यक्तियों तथा दूसरी ओर !!{structure}s
को जोड़ने वाली मूल धारणाएँ किसी पद के \emph{!!{value}} और किसी
!!a{formula} की \emph{संतुष्टि} की हैं। अनौपचारिक रूप से, किसी पद का
!!{value}, किसी !!a{structure} का !!a{element} होता है---यदि पद केवल कोई
अचर है, तो उसका !!{value} वह वस्तु है जिसे !!{structure} उस अचर को
नियत करती है; और यदि पद !!{function}s का उपयोग करके बनाया गया है, तो उसका
  !!{value} अचरों के !!{value}s और पद में आए फलन-चिह्नों को नियत किए
गए फलनों से निकाला जाता है। कोई !!^a{formula}, किसी !!a{structure} में
  \emph{संतुष्ट} है यदि विधेयों को दिया गया निर्वचन उस !!{formula} को
  !!{structure} के प्रांत में सत्य बनाता है। संतुष्टि की यह धारणा आगमनात्मक रूप से निर्दिष्ट
की जाती है: !!{structure} का विनिर्देशन सीधे बताता है कि परमाण्विक
!!{formula} कब संतुष्ट हैं; और किसी जटिल !!{formula} की संतुष्टि उसके मुख्य
संयोजक अथवा परिमाणक तथा इस बात के आधार पर परिभाषित की जाती है कि उसके
तात्कालिक !!{subformula} संतुष्ट हैं या नहीं।

यहाँ परिमाणकों का प्रकरण कुछ कठिन है, क्योंकि किसी परिमाणित !!{formula} के
तात्कालिक !!{subformula} में कोई मुक्त !!{variable} होता है, और
!!{structure}s, !!{variable}s के !!{value} निर्दिष्ट नहीं करतीं। इस कठिनाई
से निपटने के लिए हम \emph{चर-नियतन} भी प्रस्तुत करते हैं और संतुष्टि को केवल
!!a{structure} के सापेक्ष नहीं, बल्कि !!a{structure} के साथ !!a{variable}
नियतन के सापेक्ष परिभाषित करते हैं।
\end{explain}

\begin{defn}[चर-नियतन]
\emph{चर-नियतन}~$s$, किसी !!a{structure}~$\Struct{M}$ के लिए ऐसा फलन है
जो प्रत्येक !!{variable} को~$\Domain M$ के किसी !!a{element} पर
प्रतिचित्रित करता है, अर्थात् $s\colon \Var \to \Domain M$।
\end{defn}

\begin{explain}
!!^a{structure} प्रत्येक !!{constant} को कोई !!a{value}, और चर-नियतन
प्रत्येक चर को कोई मान देता है। किंतु हम उनसे बने पदों का उपयोग करके
!!{domain} में मौजूद प्रत्येक !!{element} का नामकरण भी करना चाहते हैं। इसके लिए हम पदों के
!!{value} को आगमनात्मक रूप से परिभाषित करते हैं। !!{constant}s और चरों के
लिए मान वही है जो !!{structure} अथवा चर-नियतन निर्दिष्ट करता है; अधिक जटिल
पदों के लिए उसे उन फलनों का उपयोग करके पुनरावर्ती रूप से निकाला जाता है
जिन्हें !!{structure}, !!{function}s को नियत करती है।
\end{explain}

\begin{defn}[पदों के !!^{value}]
यदि~$t$, भाषा~$\Lang L$ का कोई पद है,~$\Struct M$,~$\Lang L$ के लिए कोई
!!{structure} है, और~$s$,~$\Struct M$ के लिए कोई !!a{variable} नियतन है,
तो \emph{!!{value}}~$\Value{t}{M}[s]$ निम्न प्रकार परिभाषित है:
\begin{enumerate}
\item \indcase{t}{c}{$\Value{\indfrm}{M}[s] = \Assign{\indcomplex}{M}$।}
\item \indcase{t}{x}{$\Value{\indfrm}{M}[s] = s(\indcomplex)$।}
\item \indcase{t}{\Atom{f}{t_1, \ldots, t_n}}{
\[
\Value{\indfrm}{M}[s] = \Assign{f}{M}(\Value{t_1}{M}[s], \ldots,
\Value{t_n}{M}[s]).
\]}
\end{enumerate}
\end{defn}

\begin{defn}[$x$-परिवर्ती]
यदि~$s$, किसी !!a{structure}~$\Struct M$ के लिए कोई !!a{variable} नियतन
है, तो ऐसा कोई भी !!{variable} नियतन~$s'$,~$\Struct M$ के लिए, जो~$s$ से
अधिक-से-अधिक इस बात में भिन्न हो कि वह~$x$ को क्या नियत करता है,
वह \emph{$x$-परिवर्ती} कहलाता है---नियतन~$s$ का। यदि~$s'$ कोई $x$-परिवर्ती है
जो~$s$ का है तो हम
$\varAssign{s'}{s}{x}$ लिखते हैं।
\end{defn}

\begin{explain}
ध्यान दें कि कोई $x$-परिवर्ती, नियतन~$s$ से संबद्ध होते हुए भी,~$x$ को कोई भिन्न वस्तु
नियत करना \emph{आवश्यक} नहीं है। वास्तव में, प्रत्येक नियतन स्वयं अपना
$x$-परिवर्ती माना जाता है।
\end{explain}

\begin{defn}
  यदि~$s$, किसी !!a{structure}~$\Struct M$ के लिए कोई !!a{variable} नियतन
  है और $m \in \Domain{M}$, तो नियतन~$\Subst{s}{m}{x}$ वह चर-नियतन है
  जिसे इस प्रकार परिभाषित किया जाता है:
  \[\Subst{s}{m}{x}(y) = \begin{cases}
    m & \text{यदि } y \ident x\\
    s(y) & \text{अन्यथा}.
  \end{cases}\]
\end{defn}

दूसरे शब्दों में,~$\Subst{s}{m}{x}$ वह विशेष $x$-परिवर्ती है जो~$s$ का है,
प्रांत के !!{element}~$m$ को~$x$ का मान बनाता है, और~$x$ से भिन्न
!!{variable}s को वही वस्तुएँ नियत करता है जो~$s$ करता है।

\begin{defn}[संतुष्टि]
\ollabel{defn:satisfaction}
किसी !!a{formula}~$!A$ की किसी !!a{structure}~$\Struct M$ में किसी
!!a{variable} नियतन~$s$ के सापेक्ष संतुष्टि, चिह्नों में:
$\Sat{M}{!A}[s]$, निम्न प्रकार पुनरावर्ती रूप से परिभाषित है।
($\Sat/{M}{!A}[s]$ से हमारा आशय ``$\Sat{M}{!A}[s]$ नहीं'' है।)
\begin{enumerate}
\tagitem{prvFalse}{%
  \indcase{!A}{\lfalse}{$\Sat/{M}{\indfrm}[s]$।}}{}

\tagitem{prvTrue}{%
  \indcase{!A}{\ltrue}{$\Sat{M}{\indfrm}[s]$।}}{}

\item \indcase{!A}{\Atom{R}{t_1, \dots, t_n}}{$\Sat{M}{\indfrm}[s]$
  ठीक तब, जब $\langle \Value{t_1}{M}[s], \dots, \Value{t_n}{M}[s] \rangle \in
  \Assign{R}{M}$।}

\item \indcase{!A}{\eq[t_1][t_2]}{$\Sat{M}{\indfrm}[s]$ तभी और केवल तभी
  जब $\Value{t_1}{M}[s] = \Value{t_2}{M}[s]$।}

\tagitem{prvNot}{%
  \indcase{!A}{\lnot !B}{$\Sat{M}{\indfrm}[s]$ तभी और केवल तभी जब
    $\Sat/{M}{!B}[s]$।}}{}

\tagitem{prvAnd}{%
  \indcase{!A}{(!B \land !C)}{$\Sat{M}{\indfrm}[s]$ तभी और केवल तभी जब
    $\Sat{M}{!B}[s]$ और $\Sat{M}{!C}[s]$।}}{}

\tagitem{prvOr}{%
  \indcase{!A}{(!B \lor !C)}{$\Sat{M}{\indfrm}[s]$ तभी और केवल तभी जब
    $\Sat{M}{!B}[s]$ अथवा $\Sat{M}{!C}[s]$ (अथवा दोनों)।}}{}

\tagitem{prvIf}{%
  \indcase{!A}{(!B \lif !C)}{$\Sat{M}{\indfrm}[s]$ तभी और केवल तभी जब
    $\Sat/{M}{!B}[s]$ अथवा $\Sat{M}{!C}[s]$ (अथवा दोनों)।}}{}

\tagitem{prvIff}{%
  \indcase{!A}{(!B \liff !C)}{$\Sat{M}{\indfrm}[s]$ तभी और केवल तभी जब
    या तो $\Sat{M}{!B}[s]$ और $\Sat{M}{!C}[s]$ दोनों, अथवा न
    $\Sat{M}{!B}[s]$ और न $\Sat{M}{!C}[s]$।}}{}

\tagitem{prvAll}{%
  \indcase{!A}{\lforall[x][!B]}{$\Sat{M}{\indfrm}[s]$ तभी और केवल तभी जब
    प्रत्येक !!{element}~$m \in \Domain M$ के लिए
    $\Sat{M}{!B}[\Subst{s}{m}{x}]$।}}{}

\tagitem{prvEx}{%
  \indcase{!A}{\lexists[x][!B]}{$\Sat{M}{\indfrm}[s]$ तभी और केवल तभी जब
  कम-से-कम एक !!{element}~$m \in \Domain M$ के लिए
  $\Sat{M}{!B}[\Subst{s}{m}{x}]$।}}{}
\end{enumerate}
\end{defn}

\begin{explain}
अंतिम \iftag{notprvEx,notprvAll}{खंड में}{दो खंडों में} चर-नियतन महत्त्वपूर्ण
हैं।\iftag{prvAll}{ हम $\lforall[x][!B(x)]$ की संतुष्टि को ``सभी
$m \in
\Domain{M}$ के लिए $\Sat{M}{!B(m)}$'' द्वारा परिभाषित नहीं कर
सकते।}{}\iftag{prvEx}{ हम $\lexists[x][!B(x)]$ की संतुष्टि को ``कम-से-कम
एक $m \in
\Domain{M}$ के लिए $\Sat{M}{!B(m)}$'' द्वारा परिभाषित नहीं कर
सकते।}{} कारण यह है कि यदि $m \in
\Domain M$, तो वह भाषा का चिह्न नहीं
है, और इसलिए $!B(m)$ कोई !!a{formula} नहीं है (अर्थात्
$\Subst{!B}{m}{x}$ अपरिभाषित है)। हम यह भी नहीं मान सकते कि हमारे पास
~$\Struct{M}$ के प्रत्येक !!a{element} का नाम रखने वाले !!{constant} अथवा
पद उपलब्ध हैं, क्योंकि !!{structure} की परिभाषा में ऐसी कोई अपेक्षा नहीं
है। मानक भाषा में !!{constant}s का समुच्चय !!{denumerable} है; अतः यदि
$\Domain{M}$ !!{enumerable} नहीं है, तो प्रत्येक वस्तु का नाम रखने के लिए
पर्याप्त !!{constant} भी नहीं हैं।

हम !!{variable} नियतन प्रस्तुत करके इस समस्या को हल करते हैं; वे हमें चरों को प्रांत
के प्रत्येक !!{element} से सीधे जोड़ने देते हैं। तब, उदाहरण के लिए, यह कहने के
स्थान पर कि \iftag{prvEx}{$\lexists[x][!B(x)]$}{$\lforall[x][!B(x)]$},
~$\Struct M$ में तभी और केवल तभी संतुष्ट है जब
\iftag{prvEx}{कम-से-कम एक $m \in
\Domain{M}$ के लिए}{सभी
$m \in \Domain{M}$ के लिए $\Sat{M}{!B(m)}$}, हम कहते हैं कि वह
~$\Struct M$ में~$s$ के \emph{सापेक्ष} तभी और केवल तभी संतुष्ट है जब
$!B(x)$,~$\Subst{s}{m}{x}$ के सापेक्ष
\iftag{prvEx}{कम-से-कम एक}{प्रत्येक} $m \in \Domain M$ के लिए संतुष्ट है।
\end{explain}

\begin{ex}
मान लें कि $\Lang{L} = \{a, b, f, R\}$, जहाँ~$a$ और~$b$ !!{constant},
~$f$ कोई दो-स्थानी !!{function}, और~$R$ कोई दो-स्थानी !!{predicate} है।
निम्न प्रकार परिभाषित !!{structure}~$\Struct{M}$ पर विचार करें:
\begin{enumerate}
\item $\Domain M = \{1, 2, 3, 4\}$
\item $\Assign{a}{M} = 1$
\item $\Assign{b}{M} = 2$
\item $\Assign{f}{M}(x, y) = x+y$ यदि $x+y \le 3$, और अन्यथा $= 3$।
\item $\Assign{R}{M} = \{\tuple{1, 1}, \tuple{1, 2}, \tuple{2, 3}, \tuple{2, 4}\}$
\end{enumerate}
फलन $s(x) = 1$, जो प्रत्येक !!{variable} को $1 \in \Domain{M}$ नियत करता
है,~$\Struct{M}$ के लिए चर-नियतन है।

तब
\begin{align*}
\Value{f(a,b)}{M}[s] & = \Assign{f}{M}(\Value{a}{M}[s], \Value{b}{M}[s]).
\intertext{चूँकि~$a$ और~$b$ !!{constant} हैं, इसलिए $\Value{a}{M}[s]
  = \Assign{a}{M} = 1$ और $\Value{b}{M}[s] = \Assign{b}{M} = 2$। अतः}
\Value{f(a,b)}{M}[s] & = \Assign{f}{M}(1, 2) = 1+2 = 3.
\intertext{$f(f(a,b),a)$ का मान निकालने के लिए हमें यह देखना होगा:}
\Value{f(f(a,b),a)}{M}[s] & = \Assign{f}{M}(\Value{f(a, b)}{M}[s],
\Value{a}{M}[s]) = \Assign{f}{M}(3, 1) = 3,
\intertext{क्योंकि $3+1 > 3$। चूँकि $s(x) = 1$ और $\Value{x}{M}[s] =
  s(x)$, हमारे पास यह भी है:}
\Value{f(f(a,b),x)}{M}[s] & = \Assign{f}{M}(\Value{f(a, b)}{M}[s],
\Value{x}{M}[s]) = \Assign{f}{M}(3, 1) = 3,
\end{align*}

किसी परमाण्विक !!{formula}~$R(t_1, t_2)$ की संतुष्टि के लिए यह आवश्यक और
पर्याप्त है कि उसके तर्कों के मानों का क्रमित युग्म $\tuple{\Value{t_1}{M}[s],
  \Value{t_2}{M}[s]}$,~$\Assign{R}{M}$ का !!a{element} हो। अतः, उदाहरण
के लिए, हमारे पास $\Sat{M}{R(b,f(a,b))}[s]$ है, क्योंकि
$\tuple{\Value{b}{M},
  \Value{f(a,b)}{M}} = \tuple{2, 3} \in \Assign{R}{M}$; किंतु
$\Sat/{M}{R(x, f(a,b))}[s]$, क्योंकि
$\tuple{1, 3} \notin \Assign{R}{M}[s]$।

यह निर्धारित करने के लिए कि कोई अपरमाण्विक सूत्र~$!A$ संतुष्ट है या नहीं,
आप आगमनात्मक परिभाषा का वह खंड लगाते हैं जो मुख्य संयोजक पर लागू होता है।
उदाहरणतः, $R(a, a) \lif (R(b,
x) \lor R(x, b))$ का मुख्य संयोजक~$\lif$
है, और
\begin{align*}
  & \Sat{M}{R(a, a) \lif (R(b, x) \lor R(x, b))}[s]
    \text{ तभी और केवल तभी जब }\\
  & \qquad
  \Sat/{M}{R(a,a)}[s] \text{ अथवा } \Sat{M}{R(b, x) \lor R(x, b)}[s]
  \intertext{चूँकि $\Sat{M}{R(a,a)}[s]$ (क्योंकि $\tuple{1,1} \in
    \Assign{R}{M}$), हम अभी उत्तर निर्धारित नहीं कर सकते और पहले यह पता
    करना होगा कि $\Sat{M}{R(b, x) \lor R(x, b)}[s]$ है या नहीं:}
  & \Sat{M}{R(b, x) \lor R(x, b)}[s] \text{ तभी और केवल तभी जब }\\
  & \qquad \Sat{M}{R(b,x)}[s] \text{ अथवा } \Sat{M}{R(x, b)}[s]
  \intertext{और ऐसा है, क्योंकि $\Sat{M}{R(x, b)}[s]$
    (क्योंकि $\tuple{1,2} \in \Assign{R}{M}$)।}
\end{align*}

याद करें कि कोई $x$-परिवर्ती~$s$ का ऐसा चर-नियतन है जो~$s$ से अधिकतम इस
बात में भिन्न होता है कि वह~$x$ को क्या नियत करता है।~$\Domain{M}$ के
प्रत्येक !!{element} के लिए एक $x$-परिवर्ती~$s$ का है:
\begin{align*}
  s_1 & = \Subst{s}{1}{x}, &
  s_2 & = \Subst{s}{2}{x},\\
  s_3 & = \Subst{s}{3}{x}, &
  s_4 & = \Subst{s}{4}{x}.
\end{align*}
अतः, उदाहरण के लिए, $s_2(x) = 2$ और $s_2(y) = s(y) = 1$ उन सभी चरों~$y$
के लिए जो~$x$ से भिन्न हैं। ये सभी $x$-परिवर्ती~$s$ के हैं,
~$\Struct{M}$ के लिए, क्योंकि $\Domain{M} = \{1, 2, 3, 4\}$। विशेषतः ध्यान
दें कि $s_1 = s$ (~$s$ सदैव स्वयं अपना $x$-परिवर्ती है)।

\iftag{prvEx}{किसी अस्तित्वपरक परिमाणित !!{formula}
  ~$\lexists[x][!A(x)]$ की संतुष्टि निर्धारित करने के लिए हमें यह निर्धारित
  करना होता है कि $\Sat{M}{!A(x)}[\Subst{s}{m}{x}]$ कम-से-कम एक
  $m \in \Domain
  M$ के लिए है या नहीं। अतः
  \[
  \Sat{M}{\lexists[x][(R(b,x) \lor R(x,b))]}[s],
  \]
  क्योंकि $\Sat{M}{R(b,x) \lor R(x, b)}[\Subst{s}{1}{x}]$
  ($\Subst{s}{3}{x}$ से भी काम बनता)। किंतु
  \[
  \Sat/{M}{\lexists[x][(R(b,x) \land R(x,b))]}[s]
  \]
  क्योंकि हम चाहे जो $m \in \Domain{M}$ चुनें,
  $\Sat/{M}{R(b,x) \land R(x,b)}[\Subst{s}{m}{x}]$।}{}

\iftag{prvAll}{किसी सार्ववाची परिमाणित !!{formula}
  ~$\lforall[x][!A(x)]$ की संतुष्टि निर्धारित करने के लिए हमें यह निर्धारित
  करना होता है कि $\Sat{M}{!A(x)}[\Subst{s}{m}{x}]$ सभी $m \in \Domain M$
  के लिए है या नहीं। अतः
  \[
  \Sat{M}{\lforall[x][(R(x,a) \lif R(a,x))]}[s],
  \]
  क्योंकि $\Sat{M}{R(x,a) \lif R(a,x)}[\Subst{s}{m}{x}]$ सभी $m \in
  \Domain M$ के लिए है। $m = 1$ के लिए हमारे
  पास $\Sat{M}{R(a,x)}[\Subst{s}{1}{x}]$ है, अतः उत्तरपक्ष सत्य है;
  $m = 2$, $3$, और~$4$ के लिए हमारे पास
  $\Sat/{M}{R(x,a)}[\Subst{s}{m}{x}]$ है, अतः पूर्वपक्ष असत्य है। किंतु
  \[
  \Sat/{M}{\lforall[x][(R(a,x) \lif R(x,a))]}[s]
  \]
  क्योंकि $\Sat/{M}{R(a,x) \lif R(x,a)}[\Subst{s}{2}{x}]$ (क्योंकि
  $\Sat{M}{R(a, x)}[\Subst{s}{2}{x}]$ और $\Sat/{M}{R(x,
  a)}[\Subst{s}{2}{x}]$)।}{}

\iftag{defEx}{किसी अस्तित्वपरक परिमाणित !!{formula}
  ~$\lexists[x][!A(x)]$ की संतुष्टि निर्धारित करने के लिए हमें यह निर्धारित
  करना होता है कि $\Sat{M}{\lnot \lforall[x][\lnot !A(x)]}[s]$ है या
  नहीं। उदाहरणतः हमारे पास
  \[
  \Sat{M}{\lexists[x][(R(b,x) \lor R(x,b))]}[s].
  \]
  पहले, $\Sat{M}{R(b,x) \lor R(x, b)}[\Subst{s}{1}{x}]$
  ($\Subst{s}{3}{x}$ से भी काम बनता)। अतः
  $\Sat/{M}{\lnot(R(b,x) \lor R(x,b))}[\Subst{s}{1}{x}]$, इसलिए
  $\Sat/{M}{\lforall[x][\lnot((R(b,x) \lor R(x,b))]}[s]$, और अतः
  $\Sat{M}{\lnot\lforall[x][\lnot((R(b,x) \lor
  R(x,b))]}[s]$। दूसरी ओर,
  \[
  \Sat/{M}{\lexists[x][(R(b,x) \land R(x,b))],}[s].
  \]
  ऐसा इसलिए कि $\Sat{M}{\lforall[x][\lnot(R(b,x) \land
  R(x,b))]}[s]$, क्योंकि किसी भी $m \in \Domain M$ के लिए
  $\Sat{M}{R(b,x) \land
  R(x,b)}[\Subst{s}{m}{x}]$ नहीं। जैसा कि इन
  उदाहरणों से अनुमान लगाया जा सकता है,
  $\Sat{M}{\lexists[x][!A(x)]}[s]$ तभी और केवल तभी जब
  $\Sat{M}{!A(x)}[\Subst{s}{m}{x}]$ कम-से-कम एक $m \in \Domain
  M$ के लिए।}{}

\iftag{defAll}{किसी सार्ववाची परिमाणित !!{formula}
  ~$\lforall[x][!A(x)]$ की संतुष्टि निर्धारित करने के लिए हमें यह निर्धारित
  करना होता है कि $\Sat{M}{\lnot\lexists[x][\lnot !A(x)]}[s]$ है या नहीं।
  उदाहरणतः
  \[
  \Sat{M}{\lforall[x][(R(x,a) \lif R(a,x))]}[s],
  \]
  पहले, $\Sat{M}{R(x,a) \lif R(a,x)}[\Subst{s}{m}{x}]$ सभी
  $m \in
  \Domain M$ के लिए
  ($\Sat{M}{R(a,x)}[\Subst{s}{1}{x}]$ और
  $\Sat/{M}{R(a,x)}[\Subst{s}{m}{x}]$, जब $m = 2$, $3$, अथवा~$4$)। अतः
  कोई $m \in \Domain M$ ऐसा
  नहीं कि $\Sat{M}{\lnot(R(x,a) \lif R(a,x))}[\Subst{s}{m}{x}]$, और इसलिए
  $\Sat/{M}{\lexists[x][\lnot(R(x,a) \lif R(a,x))]}[s]$। अतः
  $\Sat{M}{\lnot\lexists[x][\lnot(R(x,a) \lif R(a,x))]}[s]$। दूसरी ओर,
  \[
  \Sat/{M}{\lforall[x][(R(a,x) \lif R(x,a))]}[s],
  \]
  क्योंकि $\Sat/{M}{R(a,x) \lif R(x,a)}[\Subst{s}{2}{x}]$, और इसलिए
  $\Sat{M}{\lexists[x][\lnot(R(a,x) \lif R(x,a))]}[s]$। जैसा कि इन
  उदाहरणों से अनुमान लगाया जा सकता है, $\Sat{M}{\lforall[x][!A(x)]}[s]$
  तभी और केवल तभी जब $\Sat{M}{!A(x)}[\Subst{s}{m}{x}]$ प्रत्येक
  $m \in \Domain
  M$ के लिए।}{}

अधिक जटिल प्रकरण के लिए इस पर विचार करें:
\[
\lforall[x][(R(a,x) \lif \lexists[y][R(x,y)])].
\]
चूँकि $\Sat/{M}{R(a,x)}[\Subst{s}{3}{x}]$ और
$\Sat/{M}{R(a,x)}[\Subst{s}{4}{x}]$, वे रोचक प्रकरण, जिनमें हमें सशर्त
के उत्तरपक्ष की चिंता करनी है, केवल $m = 1$ और $ = 2$ हैं। क्या
$\Sat{M}{\lexists[y][R(x,y)]}[\Subst{s}{1}{x}]$ सत्य है? ऐसा तब है जब
कम-से-कम एक $n \in \Domain M$ ऐसा हो कि
$\Sat{M}{R(x,y)}[\Subst{\Subst{s}{1}{x}}{n}{y}]$। वास्तव में, यदि हम
$n = 1$ लें, तो $\Subst{\Subst{s}{1}{x}}{n}{y} = \Subst{s}{1}{y} =
s$।
चूँकि $s(x) = 1$, $s(y) = 1$, और $\tuple{1,1} \in \Assign{R}{M}$, उत्तर
हाँ है।

यह निर्धारित करने के लिए कि
$\Sat{M}{\lexists[y][R(x,y)]}[\Subst{s}{2}{x}]$ है या नहीं, हमें
!!{variable} नियतनों $\Subst{\Subst{s}{2}{x}}{n}{y}$ को देखना होगा। यहाँ
$n = 1$ के लिए यह नियतन~$s_2 = \Subst{s}{2}{x}$ है, जो~$R(x,y)$ को
संतुष्ट नहीं करता ($s_2(x) =
2$, $s_2(y) = 1$, और
$\tuple{2,1}\notin \Assign{R}{M}$)। किंतु
$\Subst{\Subst{s}{2}{x}}{3}{y} = \Subst{s_2}{3}{y}$ पर विचार करें।
$\Sat{M}{R(x,y)}[\Subst{s_2}{3}{y}]$, क्योंकि $\tuple{2,3} \in
\Assign{R}{M}$; और अतः
$\Sat{M}{\lexists[y][R(x,y)]}[s_2]$।

अतः सभी $n \in \Domain M$ के लिए या तो
$\Sat/{M}{R(a,x)}[\Subst{s}{m}{x}]$ (यदि $m = 3$, $4$), अथवा
$\Sat{M}{\lexists[y][R(x,y)]}[\Subst{s}{m}{x}]$ (यदि $m = 1$, $2$), और इसलिए
\[
\Sat{M}{\lforall[x][(R(a,x) \lif \lexists[y][R(x,y)])]}[s].
\]
दूसरी ओर,
\[
\Sat/{M}{\lexists[x][(R(a,x) \land \lforall[y][R(x,y)])]}[s].
\]
हमारे पास $\Sat{M}{R(a,x)}[\Subst{s}{m}{x}]$ केवल $m = 1$ और $m =
2$
के लिए है। किंतु~$m$ के इन दोनों मानों के लिए क्रमशः कोई
$n \in
\Domain M$, अर्थात् $n = 4$, ऐसा है कि
$\Sat/{M}{R(x,y)}[\Subst{\Subst{s}{m}{x}}{n}{y}]$, और अतः
$\Sat/{M}{\lforall[y][R(x,y)]}[\Subst{s}{m}{x}]$ होता है—यह $m = 1$ और $m =
2$ दोनों के लिए। संक्षेप
में, कोई $m \in \Domain M$ ऐसा नहीं कि $\Sat{M}{R(a,x)
\land \lforall[y][R(x,y)]}[\Subst{s}{m}{x}]$।
\end{ex}

\iftag{defEx}{%
\begin{prop}\ollabel{prop:sat-ex}
$\Sat{M}{\lexists[x][!B(x)]}[s]$ तभी और केवल तभी जब कोई $x$-परिवर्ती
$s'$ ऐसा हो जो~$s$ का हो और $\Sat{M}{!B(x)}[s']$।
\end{prop}

\begin{proof}
  अभ्यास।
\end{proof}
}{}

\tagprob{defEx}
\begin{prob}
  \olref[fol][syn][sat]{prop:sat-ex} सिद्ध कीजिए।
\end{prob}
\tagendprob

\iftag{defAll}{%
\begin{prop}\ollabel{prop:sat-all}
$\Sat{M}{\lforall[x][!B(x)]}[s]$ तभी और केवल तभी जब प्रत्येक $x$-परिवर्ती
~$s'$ के लिए, जो~$s$ का है, $\Sat{M}{!B(x)}[s']$।
\end{prop}

\begin{proof}
  अभ्यास।
\end{proof}
}{}

\tagprob{defAll}
\begin{prob}
  \olref[fol][syn][sat]{prop:sat-all} सिद्ध कीजिए।
\end{prob}
\tagendprob

\begin{prob}
मान लें कि $\Lang L = \{c, f, A\}$, जिसमें एक !!{constant}, एक एक-स्थानी
!!{function} और एक दो-स्थानी !!{predicate} है; और
!!{structure}~$\Struct{M}$ इस प्रकार दी गई है:
\begin{enumerate}
\item $\Domain M = \{1, 2, 3\}$
\item $\Assign{c}{M} = 3$
\item $\Assign{f}{M}(1) = 2, \Assign{f}{M}(2) = 3, \Assign{f}{M}(3) = 2$
\item $\Assign{A}{M} = \{\tuple{1, 2}, \tuple{2, 3}, \tuple{3, 3}\}$
\end{enumerate}
(a) $s(v) = 1$ रखें, सभी !!{variable}~$v$ के लिए। पता लगाइए कि क्या
\[
\Sat{M}{\lexists[x][(A(f(z), c) \lif \lforall[y][(A(y, x) \lor A(f(y),
      x))])]}[s]
\]
है। क्यों अथवा क्यों नहीं, स्पष्ट कीजिए।

(b) ऐसी कोई भिन्न संरचना और !!{variable} नियतन दीजिए जिसमें यह
!!{formula} संतुष्ट न हो।
\end{prob}

\end{document}
